\documentclass[10pt,a4paper]{amsart}
\usepackage[margin=19mm]{geometry}
\usepackage{amsmath,amssymb,mathtools}
\usepackage[scr=rsfs,cal=euler]{mathalfa}
\usepackage{xfrac}
\usepackage[unicode,hidelinks,bookmarks=false]{hyperref}
\newcommand{\ie}{i.e.\ }
\newcommand{\cf}{cf.\ }
\newcommand{\resp}{respectively\ }
\newcommand{\Subsection}{\subsection}
\providecommand{\nobreakdash}{\nobreak\hbox{-}}
\setlength{\emergencystretch}{2em}
\multlinegap0pt
\usepackage{amssymb}
\usepackage{mathtools}
\usepackage{algorithm}\usepackage{algpseudocode}
\newtheorem{theorem}{Theorem}
\DeclareMathOperator{\supp}{supp}
\title{Linear KL-Optimal Frequency Normalisation}
\author{Kamila Szewczyk}
\date{Source: arXiv:2605.00579v2 (CC BY 4.0)}
\begin{document}
\maketitle

\noindent\emph{Source and licence.} Linear KL-Optimal Frequency Normalisation. Version arXiv:2605.00579v2 (CC BY 4.0), distributed by arXiv under the Creative Commons Attribution 4.0 International licence, http://creativecommons.org/licenses/by/4.0/, as recorded in the arXiv licence metadata for this exact version and shown on the abstract page https://arxiv.org/abs/2605.00579v2. Full licence notice: \texttt{licenses/arxiv-2605.00579-CC-BY-4.0.txt}.\par\smallskip

\noindent\textit{Editorial scope.} Counted unit: the article's theorem thm:redundancy (source0.tex lines 205--218). The article's complete original proof of it is delivered with it (source0.tex lines 220--248, one \texttt{proof} environment of 29 lines). No mathematics is written, repaired or re-derived: the statement and the proof are the article's own lines, in the article's own notation, and no retained line is substituted or re-flowed.  No further source-local context is needed: every cross-reference of every retained line resolves inside the two delivered spans.  Nothing was omitted from any retained span, so the package carries no omission record.  No retained line contains a citation, so the wrapper carries no bibliography and no bibliography entry.  No retained line contains a figure environment, an image-inclusion command or drawing code, so no image is delivered and the package declares no asset dependency.  Editorial formatting only, in addition: an AMS \texttt{amsart} preamble that restates the article's own declarations and macros used by the retained lines, namely the environments \texttt{theorem} on separate counters, together with the standard packages those lines need.  The retained environments print in this document as Theorem 1 in order of appearance, whereas the article prints the same items with the numbering of its own full document.  The complete native source of the article as distributed in the arXiv e-print for this exact version is bundled unchanged as \texttt{sources/199575/source0.tex}.
\begin{theorem}
\label{thm:redundancy}
For all $r\ge1$ and $M\ge r$,
\begin{equation}
\begin{aligned}
   \sup_{p:\,|\supp \mathcal X|=r}\rho(p,M)
   &=\log\frac{M}{M-r+1} \\
   &=\frac{r-1}{M}+\mathcal{O}\!\left((r/M)^2\right)\ \text{nats},
\end{aligned}
\end{equation}
and the supremum is approached but not attained. The same value results when $m$
is allowed to range over the reals subject to $q_a\ge1/M$, so integrality costs
nothing in the worst case.
\end{theorem}

\begin{proof}
\emph{Upper bound.} For any $p$ the allocation $x_a=1+p_a(M-r)$ is feasible
($x_a\ge1$, $\sum_{a\in \supp \mathcal X}x_a=M$) and
$x_a=(1-p_a)+p_a(M-r+1)\ge p_a(M-r+1)$, so $q_a=x_a/M\ge p_a(M-r+1)/M$ and
$p_a/q_a\le M/(M-r+1)$; hence $D(p\,\|\,q)\le\log\frac{M}{M-r+1}$. Rounding
to $m_a=\lceil p_a(M-r+1)\rceil$ keeps $m_a\ge p_a(M-r+1)$ and
$\sum_a m_a<M+1$, hence $\sum_a m_a\le M$; assigning the remaining units to any
symbol only raises its $q_a$ and so only lowers $D$, and the bound holds for
integer frequencies.
\emph{Lower bound.} Consider the ``almost hot'' distribution
$p_\varepsilon=(1-(r-1)\varepsilon,\varepsilon,\dots,\varepsilon)$, in which one
symbol carries almost all the mass, and any feasible $q$; its ``almost-cold''
mass satisfies $Q=\sum_{a\ge2}q_a\ge(r-1)/M$, since each of the $r-1$ cold
symbols needs $m_a\ge1$. The log-sum
inequality bounds the ``almost-cold'' terms below by
$(r-1)\varepsilon\log\frac{(r-1)\varepsilon}{Q}$, so with $q_1=1-Q$,
\begin{equation}
\begin{aligned}
   D(p_\varepsilon\,\|\,q)\ge
   (1-(r-1)\varepsilon)&\log\frac{1-(r-1)\varepsilon}{1-Q}\\
   +(r-1)\varepsilon&\log\frac{(r-1)\varepsilon}{Q}.
\end{aligned}
\end{equation}
As a function of $Q$ the right side is minimised at $Q=(r-1)\varepsilon$ and
increases thereafter, so for $\varepsilon<1/M$ it is increasing on the feasible
range $Q\ge(r-1)/M$ and is minimised at the endpoint $Q=(r-1)/M$, and there it tends
to $\log\frac{M}{M-r+1}$ as $\varepsilon\to0^+$. That endpoint is realised by the
integer allocation $(M-r+1,1,\dots,1)$, so the limit bounds $\rho$ as well.
\end{proof}

\end{document}
